% !TEX program = pdflatex
% Informe técnico en español (clase report, A4, 12 pt).
% Edite la portada, los capítulos en chapters/ y la bibliografía en refs.bib.
\documentclass[12pt,a4paper]{report}

%---------------------------------------------------------------------
% Idioma y codificación
%---------------------------------------------------------------------
\usepackage[utf8]{inputenc}
\usepackage[T1]{fontenc}
\usepackage[spanish,es-tabla]{babel}
\usepackage{lmodern}
\usepackage{microtype}

%---------------------------------------------------------------------
% Diseño de página y tipografía
%---------------------------------------------------------------------
\usepackage[a4paper,top=2.5cm,bottom=2.5cm,left=3cm,right=2.5cm,headheight=15pt]{geometry}
\usepackage{setspace}
\onehalfspacing
\usepackage{titlesec}
\usepackage{fancyhdr}
\usepackage[dvipsnames]{xcolor}

\definecolor{principal}{HTML}{1F3A5F}
\definecolor{acento}{HTML}{C9742A}

\titleformat{\chapter}[display]
  {\normalfont\huge\bfseries\color{principal}}
  {\filleft\large\chaptertitlename~\thechapter}
  {12pt}
  {\titlerule[1.2pt]\vspace{6pt}\filright}
\titlespacing*{\chapter}{0pt}{0pt}{30pt}
\titleformat{\section}{\normalfont\Large\bfseries\color{principal}}{\thesection}{1em}{}
\titleformat{\subsection}{\normalfont\large\bfseries\color{principal!80!black}}{\thesubsection}{1em}{}

\pagestyle{fancy}
\fancyhf{}
\fancyhead[L]{\small\nouppercase{\leftmark}}
\fancyhead[R]{\small Informe técnico}
\fancyfoot[C]{\thepage}
\renewcommand{\headrulewidth}{0.4pt}
\fancypagestyle{plain}{\fancyhf{}\fancyfoot[C]{\thepage}\renewcommand{\headrulewidth}{0pt}}

%---------------------------------------------------------------------
% Figuras, tablas, matemáticas y código
%---------------------------------------------------------------------
\usepackage{graphicx}
\usepackage{booktabs}
\usepackage{caption}
\captionsetup{font=small,labelfont={bf,color=principal}}
\usepackage{amsmath,amssymb}
\usepackage{siunitx}
\sisetup{locale=DE, output-decimal-marker={,}}
\usepackage{tikz}
\usetikzlibrary{arrows.meta,positioning,shapes.geometric,calc,backgrounds,fit}
\usepackage{pgfplots}
\pgfplotsset{compat=1.18}
\usepackage{listings}
\lstset{
  basicstyle=\ttfamily\small,
  keywordstyle=\color{principal}\bfseries,
  commentstyle=\color{gray}\itshape,
  stringstyle=\color{acento},
  numbers=left, numberstyle=\tiny\color{gray}, numbersep=8pt,
  frame=single, rulecolor=\color{gray!50}, framesep=4pt,
  breaklines=true, showstringspaces=false,
  inputencoding=utf8, extendedchars=true,
  literate={á}{{\'a}}1 {é}{{\'e}}1 {í}{{\'i}}1 {ó}{{\'o}}1 {ú}{{\'u}}1 {ñ}{{\~n}}1 {Á}{{\'A}}1 {É}{{\'E}}1 {Í}{{\'I}}1 {Ó}{{\'O}}1 {Ú}{{\'U}}1 {Ñ}{{\~N}}1,
}
\renewcommand{\lstlistingname}{Listado}
\usepackage{enumitem}

%---------------------------------------------------------------------
% Bibliografía y enlaces
%---------------------------------------------------------------------
\usepackage[round,authoryear]{natbib}
\usepackage[hidelinks,unicode]{hyperref}
\hypersetup{
  pdftitle={Compilación de documentos LaTeX en la nube},
  pdfauthor={Alex Rivera},
  pdfsubject={Informe técnico},
}

%---------------------------------------------------------------------
% Datos del documento
%---------------------------------------------------------------------
\newcommand{\titulo}{Compilación de documentos \LaTeX{} en la nube}
\newcommand{\subtitulo}{Diseño, aislamiento y evaluación de un servicio de compilación}
\newcommand{\autor}{Alex Rivera}
\newcommand{\institucion}{Universidad Nacional de Ejemplo}
\newcommand{\unidad}{Facultad de Ingeniería de Sistemas}
\newcommand{\ciudad}{Lima, Perú}
\newcommand{\fecha}{Abril de 2026}
\newcommand{\codigo}{INF-2026-014}

\begin{document}

%=====================================================================
% PORTADA
%=====================================================================
\begin{titlepage}
  \begin{tikzpicture}[remember picture,overlay]
    \fill[principal] (current page.north west) rectangle ([yshift=-5.5cm]current page.north east);
    \fill[acento] ([yshift=-5.5cm]current page.north west) rectangle ([yshift=-5.8cm]current page.north east);
    \node[anchor=north west, text=white, font=\large\bfseries, xshift=3cm, yshift=-2cm]
      at (current page.north west) {\institucion};
    \node[anchor=north west, text=white!85!principal, font=\normalsize, xshift=3cm, yshift=-2.8cm]
      at (current page.north west) {\unidad};
    \node[anchor=north east, text=white, font=\small, xshift=-2.5cm, yshift=-2cm]
      at (current page.north east) {Informe \codigo};
  \end{tikzpicture}

  \vspace*{6cm}
  \begin{flushleft}
    {\Huge\bfseries\color{principal}\titulo\par}
    \vspace{0.8cm}
    {\Large\color{black!70}\subtitulo\par}
    \vspace{1.2cm}
    {\color{acento}\rule{4cm}{2pt}\par}
    \vspace{1.2cm}
    {\large\textbf{Autor:} \autor\par}
    \vspace{0.3cm}
    {\large\textbf{Ciudad:} \ciudad\par}
    \vspace{0.3cm}
    {\large\textbf{Fecha:} \fecha\par}
  \end{flushleft}

  \vfill
  \begin{center}
    \small Este documento es una plantilla con contenido ficticio de demostración.
  \end{center}
\end{titlepage}

%=====================================================================
% RESUMEN E ÍNDICES
%=====================================================================
\pagenumbering{roman}

\chapter*{Resumen}
\addcontentsline{toc}{chapter}{Resumen}
Este informe describe el diseño y la evaluación de un servicio que compila
proyectos \LaTeX{} bajo demanda dentro de contenedores efímeros sin acceso a
la red. Se presenta la arquitectura general, el modelo de aislamiento
utilizado para ejecutar documentos de origen no confiable y una caché de
artefactos direccionada por contenido. Las mediciones realizadas sobre un
conjunto de 120 documentos muestran que, para documentos breves, el arranque
del contenedor y no el motor de composición es el principal componente de la
latencia, y que precalentar la imagen reduce la latencia mediana en un 38~\%.

\tableofcontents
\listoffigures
\listoftables

\cleardoublepage
\pagenumbering{arabic}

%=====================================================================
% CAPÍTULOS
%=====================================================================
\chapter{Introducción}
\label{cap:introduccion}

\section{Contexto}
Los sistemas de composición tipográfica basados en \TeX{} siguen siendo la
herramienta de referencia para publicar artículos científicos, tesis y
documentación técnica \citep{lamport1994}. Sin embargo, instalar y mantener
una distribución completa de \TeX{} Live exige varios gigabytes de disco y un
conocimiento considerable de su ecosistema de paquetes. Para una organización
que solo necesita transformar una plantilla y unos datos en un archivo PDF
(por ejemplo, un sistema de generación de certificados o de informes
periódicos), esa carga resulta desproporcionada.

Durante los últimos años han aparecido editores en línea que centralizan la
compilación \citep{overleaf2023}, pero sus componentes internos no se exponen
como piezas reutilizables. Este informe estudia la compilación como un
servicio independiente, sin estado, que recibe un conjunto de archivos y
devuelve el PDF resultante junto con un diagnóstico estructurado de errores y
advertencias.

\section{Objetivos}
El trabajo persigue tres objetivos concretos:

\begin{enumerate}[label=\textbf{O\arabic*.}]
  \item Diseñar un modelo de ejecución que permita compilar documentos de
    origen no confiable en una infraestructura compartida sin comprometer el
    servidor ni los datos de otros usuarios.
  \item Garantizar que una misma entrada produzca siempre la misma salida
    mientras no cambie la versión de la imagen de compilación.
  \item Reducir la latencia percibida mediante una caché de artefactos que
    evite recompilar documentos que no han cambiado.
\end{enumerate}

\section{Alcance y estructura del informe}
El alcance se limita a los tres motores más utilizados (pdf\TeX{}, Xe\TeX{} y
Lua\TeX{}) y al flujo de trabajo gestionado por \texttt{latexmk}. Quedan fuera
del estudio las funcionalidades que requieren ejecutar programas externos
durante la compilación, como \texttt{minted} o la externalización de figuras
de TikZ, porque el modelo de aislamiento las prohíbe de forma deliberada.

El resto del documento se organiza de la siguiente manera. El
capítulo~\ref{cap:diseno} describe la arquitectura del servicio, el modelo de
aislamiento y la clave de caché. El capítulo~\ref{cap:resultados} presenta la
metodología de evaluación, los resultados obtenidos y las lecciones
aprendidas durante seis meses de operación.

\chapter{Diseño del servicio}
\label{cap:diseno}

\section{Arquitectura general}
El servicio se compone de tres planos claramente separados: el plano de
entrada, que valida y normaliza los proyectos recibidos; el plano de
compilación, formado por contenedores efímeros; y el plano de almacenamiento,
donde se guardan los artefactos generados. La figura~\ref{fig:arquitectura}
muestra la relación entre ellos.

\begin{figure}[htbp]
  \centering
  \begin{tikzpicture}[
      font=\small, node distance=6mm and 10mm,
      caja/.style={draw, rounded corners=2pt, minimum width=28mm, minimum height=10mm,
                   align=center, fill=white, thick},
      flecha/.style={-{Latex[length=2.5mm]}, thick},
    ]
    \node[caja] (cliente) {Cliente\\(SDK o CLI)};
    \node[caja, right=of cliente, fill=principal!10] (api) {API};
    \node[caja, right=of api, fill=acento!15] (cache) {Caché de\\artefactos};
    \node[caja, below=of api, fill=green!10] (worker) {Trabajador};
    \node[caja, right=of worker, fill=red!8] (sandbox) {Contenedor\\aislado};
    \node[caja, below=of worker, fill=gray!10] (almacen) {Almacén\\de archivos};
    \draw[flecha] (cliente) -- node[above]{proyecto} (api);
    \draw[flecha] (api) -- node[above]{hash} (cache);
    \draw[flecha] (api) -- node[left]{fallo de caché} (worker);
    \draw[flecha] (worker) -- node[above]{lanza} (sandbox);
    \draw[flecha] (sandbox) |- node[pos=0.75, below]{PDF y registro} (almacen);
    \draw[flecha] (almacen) -- (worker);
    \draw[flecha, dashed] (cache) |- (almacen);
    \begin{scope}[on background layer]
      \node[fit=(worker)(sandbox)(almacen), draw=gray!60, dashed, rounded corners,
            inner sep=4pt, label={[gray]below:plano de compilación}] {};
    \end{scope}
  \end{tikzpicture}
  \caption{Arquitectura del servicio. El plano de compilación no tiene acceso a
    la red ni a los datos de otros usuarios.}
  \label{fig:arquitectura}
\end{figure}

Cada petición se procesa en cuatro pasos. Primero, la API comprueba los
límites de tamaño y rechaza rutas peligrosas (referencias al directorio
superior, rutas absolutas o enlaces simbólicos). Segundo, calcula un resumen
criptográfico del proyecto normalizado. Tercero, consulta la caché y, si
existe un artefacto para ese resumen, lo devuelve de inmediato. Cuarto, en
caso contrario, un trabajador lanza un contenedor, ejecuta la compilación y
almacena el resultado.

\section{Modelo de aislamiento}
Un documento \LaTeX{} es, en la práctica, un programa: puede leer y escribir
archivos, consumir memoria sin límite y, si el intérprete lo permite,
ejecutar órdenes del sistema. Por ello cada compilación se ejecuta en un
contenedor nuevo con las restricciones de la tabla~\ref{tab:limites}.

\begin{table}[htbp]
  \centering
  \caption{Límites aplicados a cada compilación.}
  \label{tab:limites}
  \begin{tabular}{@{}llr@{}}
    \toprule
    Recurso & Mecanismo & Límite \\
    \midrule
    Tiempo de pared & \texttt{timeout -s KILL} & \SI{60}{s} (síncrono) / \SI{300}{s} (asíncrono) \\
    Memoria & cgroup v2 & \SI{2}{GiB} \\
    Procesadores & cgroup v2 & 2 vCPU \\
    Procesos & \texttt{pids.max} & 64 \\
    Red & espacio de nombres & ninguna \\
    Llamadas al sistema & seccomp & lista blanca \\
    Ejecución de programas & \texttt{shell\_escape = f} & deshabilitada \\
    \bottomrule
  \end{tabular}
\end{table}

Además de los límites de recursos, el sistema de archivos raíz del contenedor
es de solo lectura y el directorio de trabajo se monta en memoria. La opción
\texttt{openout\_any = p} del archivo \texttt{texmf.cnf} impide que el motor
escriba fuera de ese directorio, y \texttt{-no-shell-escape} refuerza en la
línea de órdenes la prohibición ya establecida en la configuración.

\section{Clave de caché}
Sea $B = \{(p_i, b_i)\}_{i=1}^{n}$ el conjunto de archivos del proyecto, donde
$p_i$ es la ruta relativa normalizada y $b_i$ su contenido en bytes. Dado el
motor $e$ y la versión de la imagen $v$, la clave de caché se define como
\begin{equation}
  \kappa(B, e, v) = H\Big( e \,\|\, v \,\|\, \big\Vert_{i=1}^{n}
    \big(|p_i| \,\|\, p_i \,\|\, |b_i| \,\|\, b_i\big) \Big),
  \label{eq:clave}
\end{equation}
donde $H$ es SHA-256 y $\Vert$ denota la concatenación con prefijos de
longitud. Incluir $v$ en la clave garantiza que una actualización de
\TeX{} Live invalide los artefactos antiguos, que son la principal fuente de
resultados no reproducibles en proyectos de larga duración.

\section{Interfaz de programación}
El listado~\ref{lst:peticion} muestra una petición mínima al servicio. El
cuerpo contiene los archivos del proyecto codificados como texto y el nombre
del archivo principal; la respuesta es el PDF o, si se solicita mediante la
cabecera \texttt{Accept}, un objeto JSON con los diagnósticos.

\begin{lstlisting}[language=bash, caption={Petición de compilación con curl.}, label={lst:peticion}]
curl -X POST https://latex-api.example/v1/compile \
  -H "Authorization: Bearer $API_KEY" \
  -H "Content-Type: application/json" \
  -d '{
        "engine": "pdflatex",
        "entrypoint": "main.tex",
        "files": [
          { "path": "main.tex", "content": "\\documentclass{article}..." }
        ]
      }' \
  -o salida.pdf
\end{lstlisting}

Cuando la compilación falla, el servicio devuelve un documento de problema
conforme a RFC~9457 con el código \texttt{compile\_failed} y la lista de
diagnósticos, cada uno con su archivo, línea y mensaje. De este modo el
cliente puede señalar el error directamente en el editor sin analizar el
registro completo del motor.

\chapter{Evaluación y resultados}
\label{cap:resultados}

\section{Metodología}
Se recopilaron 120 proyectos públicos de distinto tipo (artículos de
congreso, artículos de revista, currículos y tesis) y se compiló cada uno
diez veces en una máquina virtual de dos procesadores virtuales y
\SI{4}{GiB} de memoria. Las ejecuciones se dividieron en dos escenarios:

\begin{itemize}
  \item \textbf{Frío}: contenedor nuevo, sin archivos de formato
    pregenerados ni caché de fuentes.
  \item \textbf{Caliente}: imagen con los formatos de los tres motores y la
    base de datos de fuentes de \texttt{luaotfload} generados durante la
    construcción de la imagen.
\end{itemize}

Para cada ejecución se registró el tiempo de cada fase: creación del
contenedor, carga del formato, pasadas del motor, procesamiento de la
bibliografía y escritura de los artefactos.

\section{Latencia}
La tabla~\ref{tab:latencia} resume la latencia de extremo a extremo por tipo
de documento. La creación del contenedor cuesta alrededor de \SI{1,1}{s}
independientemente del documento, y el motor solo domina el tiempo total en
documentos de más de cuarenta páginas.

\begin{table}[htbp]
  \centering
  \caption{Latencia de extremo a extremo en segundos (mediana / percentil 95).}
  \label{tab:latencia}
  \begin{tabular}{@{}lrrr@{}}
    \toprule
    Tipo de documento & Páginas & Frío & Caliente \\
    \midrule
    Artículo de congreso sin bibliografía & 4--6 & 5,8 / 7,9 & 2,9 / 3,6 \\
    Artículo de congreso con \textsc{Bib}\TeX{} & 6--8 & 9,7 / 12,4 & 6,1 / 7,8 \\
    Artículo de revista con biber & 10--14 & 14,2 / 18,0 & 10,5 / 12,9 \\
    Currículo (fontspec, Lua\TeX{}) & 1--2 & 8,3 / 10,1 & 4,0 / 5,2 \\
    Tesis & 80+ & 41,6 / 55,3 & 36,9 / 48,7 \\
    \bottomrule
  \end{tabular}
\end{table}

La figura~\ref{fig:latencia} muestra la latencia mediana en función del
número de páginas para ambos escenarios. La diferencia entre las curvas se
mantiene aproximadamente constante, lo que confirma que las optimizaciones
actúan sobre el coste fijo de arranque y no sobre el tiempo de composición.

\begin{figure}[htbp]
  \centering
  \begin{tikzpicture}
    \begin{axis}[
        width=0.85\textwidth, height=7cm,
        xlabel={páginas}, ylabel={latencia mediana (s)},
        xmin=0, xmax=60, ymin=0, ymax=30,
        grid=major, grid style={gray!25},
        legend pos=north west, legend style={font=\small},
      ]
      \addplot[color=acento, mark=square*, thick]
        coordinates {(2,6.1) (6,7.4) (10,9.0) (20,13.2) (40,21.5) (60,29.0)};
      \addplot[color=principal, mark=*, thick]
        coordinates {(2,3.6) (6,4.7) (10,6.1) (20,10.3) (40,18.4) (60,26.1)};
      \legend{frío, caliente}
    \end{axis}
  \end{tikzpicture}
  \caption{Latencia mediana de compilación según el número de páginas.}
  \label{fig:latencia}
\end{figure}

\section{Comportamiento de la caché}
Durante seis meses de operación, el 68~\% de las peticiones se resolvieron
desde la caché. La tasa de aciertos varió notablemente según el tipo de
cliente: las plataformas docentes que regeneran el PDF de una tarea en cada
visita alcanzaron el 91~\%, mientras que los editores interactivos que
compilan en cada pulsación se mantuvieron por debajo del 20~\%. La inclusión
de la versión de la imagen en la ecuación~\eqref{eq:clave} provocó un vaciado
completo de la caché en cada actualización de \TeX{} Live; la tasa de
aciertos recuperó su valor habitual en unas 36 horas.

\section{Análisis de fallos}
De las compilaciones que no estaban en caché, el 11,4~\% terminó con error.
Los diagnósticos atribuyeron el 54~\% de esos errores a archivos
inexistentes referenciados con \verb|\includegraphics| o \verb|\input|, el
23~\% a secuencias de control no definidas y el 9~\% a paquetes que requieren
la ejecución de programas externos. Solo el 0,3~\% de las compilaciones fue
interrumpido por el tiempo límite y ninguna superó los límites de memoria o
de procesos una vez ajustados los valores de la tabla~\ref{tab:limites}.

\section{Conclusiones}
Los resultados permiten extraer tres conclusiones principales:

\begin{enumerate}
  \item En documentos breves, la sobrecarga de infraestructura pesa más que
    la composición tipográfica; precalentar la imagen elimina la mayor parte
    de ese coste.
  \item Una caché direccionada por contenido resulta muy rentable porque las
    cargas de trabajo reales vuelven a solicitar documentos sin cambios con
    mucha más frecuencia de la que los modifican.
  \item El aislamiento estricto no supuso un obstáculo en la práctica: menos
    del 10~\% de los errores se debieron a funcionalidades prohibidas, y en
    todos los casos existía una alternativa compatible.
\end{enumerate}

Como trabajo futuro se propone la compilación incremental por páginas y la
transmisión de los diagnósticos mientras el motor sigue en ejecución, de
forma que el usuario reciba el primer error antes de que termine la
compilación completa \citep{knuth1986}.


%=====================================================================
% BIBLIOGRAFÍA
%=====================================================================
\cleardoublepage
\addcontentsline{toc}{chapter}{Bibliografía}
\bibliographystyle{apalike}
\bibliography{refs}

\end{document}
